\documentclass[book.tex]{subfiles}

\begin{document}

\chapter{Semantics for Array Programs}

\label{chap:grammar}
\noindent\rule{11cm}{0.4pt} \\
\textbf{Prerequisites:} \autoref{chap:autodiff}, \autoref{chap:hindley_milner}  \\
\textbf{Difficulty Level:} *** \\
\noindent\rule{11cm}{0.4pt}
\vspace{5mm}

Modern scientific programs are written to operate on arrays of numbers. These arrays provide a natural representation for numerical calculations. In this chapter, we consider a type theory for array calculations

\section{Syntax}

Physika has two types of \lstinline{for} loops. The first \lstinline{for} loop behaves as a control structure.
\begin{lstlisting}[language=python]
sum = 0
for i in 100:
  sum += i
\end{lstlisting}

We also make use of a second \lstinline{for} which acts as an array builder. These \lstinline{for} loops must feature only a single expression within the \lstinline{for} loop which acts as an array constructor.

\begin{lstlisting}[language=python]
A = for i in N: i
\end{lstlisting}

We also allow Physika to perform indexed access to array elements
\begin{lstlisting}[language=python]
A[i]
\end{lstlisting}

\section{Types and Semantics}

Typing rules for array creation and indexing are defined as follows 

\begin{alignat*}{2}
    &\inference[1.]
    {
     x: \tau_1, \quad \Gamma \vdash e : \tau_2
    }
    {
    \Gamma \vdash (\textrm{for}\ x : \tau_1. e) : \tau_1 \Rightarrow \tau_2
    } \ \textrm{[For]} &&\qquad
    \inference[2.]
    {
     \Gamma \vdash v_1: (\tau_1 \Rightarrow \tau_2), \quad v_2 : \tau_1
    }
    {
     \Gamma \vdash v_1[v_2] : \tau_2 
    } \ \textrm{[Index]} \\
\end{alignat*}

For loops which operate as control structures are also governed by these same rules. Stateful update operations in a for-loop are wrapped up as an array "effect" objects which is automatically evaluated.

\end{document}